
%%%%%%%%%%%%%%%%%%%%%%%%%%%%%%%%%%%%%%%%
%           main body
%%%%%%%%%%%%%%%%%%%%%%%%%%%%%%%%%%%%%%%%


%-----------------------------------
%	TITLE PAGE
%-----------------------------------

\title[Probability \& Statistics]{\LARGE \bfseries 第三部分\quad 概率与统计} % The short title appears at the bottom of every slide, the full title is only on the title page
\author[Exercises] % Your institution as it will appear on the bottom of every slide, may be shorthand to save space
{\Large \bfseries 习题}
% \author{Singular Value Decomposition} % Your name
% \date{\today} % Date, can be changed to a custom date
\date{}

%------------------------------------------------

\begin{document}


%------------------------------------------------

\begin{frame}

\titlepage % Print the title page as the first slide

\end{frame}

%------------------------------------------------

% \begin{frame}
% \frametitle{内容提要} % Table of contents slide, comment this block out to remove it
% \tableofcontents % Throughout your presentation, if you choose to use \section{} and \subsection{} commands, these will automatically be printed on this slide as an overview of your presentation
% \end{frame}

%-----------------------------------
%	PRESENTATION SLIDES
%-----------------------------------

%------------------------------------------------

\begin{frame}

{\color{red} \scriptsize 带$^*$号的题目可以不做}

\vspace{1em}

1. 蒲丰投针问题：设有一个以平行且等距木纹铺成的地板，现在随意抛一支长度$\ell$小于等于平行纹线之间距离$h$的针
\vspace{0.6em}
\begin{itemize}
\item[(1)] 设随机变量$X$为针的中心和最近的平行纹线的距离，$\Theta$为针所在的直线和最近的平行纹线所夹的锐角或直角。它们都是均匀分布，写出它们的概率密度函数
\vspace{0.4em}
\item[(2)] $X$与$\Theta$是独立随机变量，求二者的联合密度函数
\vspace{0.4em}
\item[(3)] $X$与$\Theta$满足什么条件时，针与纹线相交？其概率等于多少？
\vspace{0.4em}
\item[(4$^*$)] 若针的长度$\ell$大于平行纹线之间距离$h$, 针与纹线相交的概率等于多少？
\end{itemize}

\end{frame}

%------------------------------------------------

\begin{frame}

2. 设有$n$个球，等可能地被放入$N$个不同的盒子中（$n \le N$）
\vspace{0.6em}
\begin{itemize}
\item[(1)] 若小球是不同的，盒子容量不限，求指定的$n$个箱子中，恰好各有一个球的概率
\vspace{0.4em}
\item[(2)] 若小球是相同的，每个盒子至多容纳一个小球，求指定的$n$个箱子中都有小球的概率
\vspace{0.4em}
\item[(3)] 若小球是相同的，盒子容量不限，求指定的$n$个箱子中，恰好各有一个球的概率
\end{itemize}

\vspace{1em}

3. 连续不断且重复地进行参数为$p$的伯努利试验，定义随机变量$X$为试验恰好取得$r$次``成功''的结果时，试验进行的次数。这样的随机变量$X$服从负二项分布$NB(r,p)$. 请推导其概率质量函数

\end{frame}

%------------------------------------------------

\begin{frame}

4. 假设某地区某种病毒携带率为$1\%$. 现有一种检测方法，对病毒携带者有$95\%$的诊断率，但是对未携带病毒者有$1\%$的误诊率。现随机选定某人独立地检测$3$次，发现$2$次呈阳性反应，$1$次呈阴性反应的概率是多少？

\vspace{1em}

5. 假设某种电子元件的使用寿命$X$（单位为小时）服从指数分布$Exp(1/1000)$, 问这种电子元件的使用寿命超过1000小时的概率是多少

\vspace{1em}

6. 假设某地区人的智商服从正态分布$N(100, 15^2)$, 求在此地区任选一个人，其智商超过150以及200的概率分别是多少？

\vspace{1em}

7. 验证正态分布的共轭先验也是正态分布

\end{frame}

%------------------------------------------------

\begin{frame}

8. 服从拉普拉斯分布的随机变量$X$的概率密度函数为
$$f(x) = Ae^{-|x|}$$
\begin{itemize}
\item[(1)] 求常数$A$
\vspace{0.4em}
\item[(2)] 求$X$的分布函数$F(x)$
\end{itemize}

\vspace{1em}

9. 设随机变量$X_1,\ldots,X_n$独立同指数分布$Exp(\lambda)$, 求证随机变量$Y = X_1 + \cdots + X_n$的概率密度函数为
$$f(y) = \dfrac{\lambda^n y^{n-1} e^{-\lambda y}}{(n-1)!}$$

10$^*$. 求（验证）\S 4常见分布中所列举的分布的期望与方差

\end{frame}

%------------------------------------------------

\begin{frame}

11. 设随机变量$X$服从标准正态分布$N(0,1)$, 令$Y=X^2$. 求$(X,Y)$的联合密度函数，并判断$X,Y$是否线性相关

\vspace{1em}

12. 设随机向量$(X,Y)$服从二元正态分布
\begin{itemize}
\item[(1)] 请用$X,Y$的期望$\mu_1,\mu_2,$ 方差$\sigma_1,\sigma_2$以及相关系数$\rho$来写出$(X,Y)$的概率密度函数的具体表达式
\vspace{0.4em}
\item[(2)] 求$X$或$Y$的边缘密度函数
\end{itemize}

\vspace{1em}

13. 设随机向量$X$服从标准正态分布$N(0,1)$, $Y$服从Rademacher分布，即$P(Y = 1) = P(Y = -1) = \frac{1}{2}$. 令$Z = XY$, 求证
\begin{itemize}
\item[(1)] $Z$服从标准正态分布$N(0,1)$
\vspace{0.4em}
\item[(2)] $Cov(X,Z) = 0$
\vspace{0.4em}
\item[(2)] $X,Z$不独立
\end{itemize}

\end{frame}

%------------------------------------------------

\begin{frame}

14. 证明连续型随机变量的切比雪夫不等式

\vspace{1em}

15. 从切比雪夫不等式出发推导切比雪夫定理

\vspace{1em}

16.有一千万人投保了某保险公司的某项保险，该项保险规定每个投保人每年需缴纳十元保费，如果被保险人死亡，则保险受益人可向保险公司领取十万元的赔偿款。假设在一年内，被保险人的死亡概率都是十万分之五，求
\vspace{0.4em}
\begin{itemize}
\item[(1)] 第一年保险公司利润不低于五千万元的概率
\vspace{0.4em}
\item[(2)] 第一年保险公司亏本的概率
\end{itemize}

\vspace{1em}

17. 设$x_1,\ldots,x_n$是正态总体$N(\mu,\sigma^2)$的样本观测值，求未知参数$\mu,\sigma^2$的最大似然估计

\end{frame}

%------------------------------------------------

\begin{frame}

18$^*$. 在采用线性模型$Y = f(\mathbf{x}) = \sum\limits_{j=0}^m w_jx_j  = \mathbf{w}^T\mathbf{x} \quad (x_0=1)$以及样本集$\{(y_i,\mathbf{x}_i)\}_{i=1}^n$的问题中，我们假设
\begin{itemize}
\item $y_i = f(\mathbf{x}_i) = \mathbf{w}^T \mathbf{x}_i + \varepsilon_i$
\item $\varepsilon_i \sim N(0, \sigma^2)$, $w_j \sim N(0, \xi^2)$
\item $\varepsilon_1, \ldots, \varepsilon_n, w_1, \ldots, w_m$相互独立
\end{itemize}
如果用所谓的最大后验估计($f(\mathbf{w} | \{\mathbf{x}_i,y_i\}) \propto f(\{\mathbf{x}_i,y_i\} | \mathbf{w}) f(\mathbf{w})$)
$$\argmax\limits_{\mathbf{w}} \left( \sum\limits_{i=1}^n \ln f(y_i|\mathbf{w}) + \sum\limits_{j=1}^m \ln f(w_j) \right)$$
请证明这等价于岭回归（Ridge Regression）
$$\argmin\limits_{\mathbf{w}} \left( \sum\limits_{i=1}^n \ln (y_i - \mathbf{w}^T\mathbf{x}_i) + \lambda \Vert \mathbf{w} \Vert_2^2 \right)$$
\end{frame}

%------------------------------------------------

% \begin{frame}

% \begin{block}{Vandermonde行列式的计算}
% \begin{enumerate}
% % \addtocounter{enumi}{0}
% \item 把第$n-1$行乘以$-a_1$加到第$n$行，再把第$n-2$行乘以$-a_1$加到第$n-1$行。按照这种顺序依次向上，直到用第$2$行乘以$-a_1$加到第$1$行。由行列式性质，$V_n$的值不变。
% \[
% V_n = \det \begin{pmatrix}
% 1 & 1 & \cdots & 1 \\
% 0 & a_2-a_1 & \cdots & a_n-a_1 \\
% \vdots & \vdots & & \vdots \\
% 0 & a_2^{n-2}(a_2-a_1) & \cdots & a_n^{n-2}(a_n-a_1) \\
% \end{pmatrix}
% \]
% \end{enumerate}
% \end{block}

% \end{frame}

%------------------------------------------------
% % closing page
% \begin{frame}
% \HUGE{\centerline{\bfseries {\color{gray} The} {\color{zkblue} End}}}

% \pause

% \vspace{1.2em}

% \Large{\centerline{\bfseries {\color{gray}Thank} \quad {\color{zkblue} You}}}

% \end{frame}

%----------------------------------------------------------------------------------------

\end{document}
